\documentclass[12pt]{article}
\usepackage[a4paper,margin=1in]{geometry}
\usepackage{iftex}
\ifXeTeX
\usepackage{fontspec}
\defaultfontfeatures{Ligatures=TeX}
\setmainfont{Times New Roman}
\else
\usepackage[T1]{fontenc}
\usepackage[utf8]{inputenc}
\usepackage{newtxtext}
\usepackage{newtxmath}
\fi
\usepackage{enumitem}
\usepackage{hyperref}
\usepackage{parskip}
\usepackage{microtype}
\usepackage[normalem]{ulem}
\usepackage{xcolor}

\newcommand{\practiceid}[1]{\texttt{#1}}
\newlength{\readinggutter}
\setlength{\readinggutter}{2.8em}
\newlength{\readingfirstindent}
\setlength{\readingfirstindent}{1.5em}
\newlength{\questionblockgap}
\setlength{\questionblockgap}{0.45\baselineskip}
\newlength{\exampleblockgap}
\setlength{\exampleblockgap}{0.65\baselineskip}
\newlength{\passageparagraphgap}
\setlength{\passageparagraphgap}{0.45\baselineskip}
\newenvironment{exampleblock}{\par\vspace{0.35\baselineskip}\begingroup\raggedright\setlength{\parskip}{0pt}\setlength{\parindent}{0pt}}{\par\endgroup\vspace{\exampleblockgap}}
\newcommand{\readinglineindent}[2]{%
\par\noindent%
\hangindent=\dimexpr\readinggutter+0.9em\relax%
\hangafter=1%
\makebox[\readinggutter][r]{\scriptsize\color{gray}#1}\hspace{0.9em}\hspace*{\readingfirstindent}#2%
\par}
\newcommand{\readingline}[2]{%
\par\noindent%
\hangindent=\dimexpr\readinggutter+0.9em\relax%
\hangafter=1%
\makebox[\readinggutter][r]{\scriptsize\color{gray}#1}\hspace{0.9em}#2%
\par}

\setlength{\parindent}{0pt}
\setlength{\parskip}{0.5em}
\setlist[enumerate]{topsep=0.25em,itemsep=0.18em,parsep=0pt,partopsep=0pt}
\setlength{\emergencystretch}{2em}

\begin{document}
\section*{TOEFL Practice 05 - Tryout}
Source of truth: DOCX only.
\section*{Listening Comprehension}
\subsection*{Part A}
DIRECTIONS: In Part A you will hear short conversations between two speakers. At the end of each conversation, a third speaker will ask a question about what the first two speakers said. Each conversation and each question will be spoken only one time. Therefore, you must listen carefully to understand what each speaker says. After you hear a conversation and the question, read the four choices and select the one that is the best answer to the question the speaker asked. Then, on your answer. sheet, find the number of the question and blacken the space that corresponds to the letter for the answer you have chosen. Blacken the space completely so that the letter inside the space does not show.\par
\begin{exampleblock}
Listen to the following example.\par
On the recording, you hear:\par
(Man) Does the car need to be filled?\par
(Woman) Mary stopped at the gas station on her way home.\par
(Narrator) What does the woman mean?\par
In your test book, you will read:\par
(A) Mary bought some food.\par
(B) Mary had car trouble.\par
(C) Mary went shopping.\par
(D) Mary bought some gas.\par
From the conversation you learn that Mary stopped at the gas station on her way home. The best answer to the question "Does the car need to be filled?" is (D), "Mary bought some gas." Therefore, the correct answer is (D).\par
Now let us begin Part A with question number 1.\par
\end{exampleblock}
\noindent\textbf{1.}
\begin{enumerate}[label=(\Alph*),leftmargin=*,itemsep=0.2em]
\item They gave me a lift.
\item He didn't stand close to the door.
\item He walked through the door on the left.
\item The door was left open.
\end{enumerate}
\par
\vspace{0.25\baselineskip}
\noindent\textbf{2.}
\begin{enumerate}[label=(\Alph*),leftmargin=*,itemsep=0.2em]
\item Rich students don't take math.
\item Rich didn't know much math.
\item Rich is taking a math course.
\item Rich should take math next year..
\end{enumerate}
\par
\vspace{0.25\baselineskip}
\noindent\textbf{3.}
\begin{enumerate}[label=(\Alph*),leftmargin=*,itemsep=0.2em]
\item They wanted to know if they could buy the picture.
\item They were wandering around looking for a picture.
\item They stopped by the picture gallery to ask questions.
\item They thought that the picture was wonderful.
\end{enumerate}
\par
\vspace{0.25\baselineskip}
\noindent\textbf{4.}
\begin{enumerate}[label=(\Alph*),leftmargin=*,itemsep=0.2em]
\item He spent 15 minutes reading.
\item He talked for 50 minutes
\item He found the magazine.
\item He enjoyed the article.
\end{enumerate}
\par
\vspace{0.25\baselineskip}
\noindent\textbf{5.}
\begin{enumerate}[label=(\Alph*),leftmargin=*,itemsep=0.2em]
\item He's looking for a new apartment.
\item He decided to buy a new car.
\item He doesn't have the figures.
\item His apartment is better than hers.
\end{enumerate}
\par
\vspace{0.25\baselineskip}
\noindent\textbf{6.}
\begin{enumerate}[label=(\Alph*),leftmargin=*,itemsep=0.2em]
\item The police stop most drivers.
\item The speed limit is unreasonable.
\item Drivers don't watch the traffic carefully.
\item Few people drive within the speed limit.
\end{enumerate}
\par
\vspace{0.25\baselineskip}
\noindent\textbf{7.}
\begin{enumerate}[label=(\Alph*),leftmargin=*,itemsep=0.2em]
\item Pink looks great on you.
\item Black is out of style.
\item The fashion changes every year.
\item I don't know the date of the fashion show.
\end{enumerate}
\par
\vspace{0.25\baselineskip}
\noindent\textbf{8.}
\begin{enumerate}[label=(\Alph*),leftmargin=*,itemsep=0.2em]
\item The banks are open at this hour.
\item The river is overflowing its banks.
\item Their bank is located nearby.
\item The banks close too early.
\end{enumerate}
\par
\vspace{0.25\baselineskip}
\noindent\textbf{9.}
\begin{enumerate}[label=(\Alph*),leftmargin=*,itemsep=0.2em]
\item The teacher taught about dreams.
\item The students are required to take the course.
\item The lecture on dreams was a great success.
\item Hardly anyone was listening to the teacher.
\end{enumerate}
\par
\vspace{0.25\baselineskip}
\noindent\textbf{10.}
\begin{enumerate}[label=(\Alph*),leftmargin=*,itemsep=0.2em]
\item They haven't had much luck with the weather.
\item The door was locked because of severe weather.
\item She's not sure if the door is locked.
\item She can't remember where the key to the door is.
\end{enumerate}
\par
\vspace{0.25\baselineskip}
\noindent\textbf{11.}
\begin{enumerate}[label=(\Alph*),leftmargin=*,itemsep=0.2em]
\item She gave away my half.
\item She is a half hour late.
\item She is very angry.
\item She is very hungry.
\end{enumerate}
\par
\vspace{0.25\baselineskip}
\noindent\textbf{12.}
\begin{enumerate}[label=(\Alph*),leftmargin=*,itemsep=0.2em]
\item It's easy to get along with Paul.
\item Paul was not feeling well.
\item She is sorry that Paul is ill.
\item Paul was uncomfortable.
\end{enumerate}
\par
\vspace{0.25\baselineskip}
\noindent\textbf{13.}
\begin{enumerate}[label=(\Alph*),leftmargin=*,itemsep=0.2em]
\item They showed us the back entrance.
\item When they return, the show will have started.
\item They don't know what time the show starts.
\item When they came back, there was no time left.
\end{enumerate}
\par
\vspace{0.25\baselineskip}
\noindent\textbf{14.}
\begin{enumerate}[label=(\Alph*),leftmargin=*,itemsep=0.2em]
\item She doesn't need any help.
\item She has two people helping her.
\item She's had enough.
\item The man is very kind.
\end{enumerate}
\par
\vspace{0.25\baselineskip}
\noindent\textbf{15.}
\begin{enumerate}[label=(\Alph*),leftmargin=*,itemsep=0.2em]
\item There is no light in number 8.
\item Smoking in the lounge is prohibited.
\item The management doesn't allow complaints.
\item The conversation in the lounge is too loud.
\end{enumerate}
\par
\vspace{0.25\baselineskip}
\noindent\textbf{16.}
\begin{enumerate}[label=(\Alph*),leftmargin=*,itemsep=0.2em]
\item Six copies were made last week.
\item The copier doesn't work well.
\item Eleven copies is plenty.
\item There are 11 jars.of jam.
\end{enumerate}
\par
\vspace{0.25\baselineskip}
\noindent\textbf{17.}
\begin{enumerate}[label=(\Alph*),leftmargin=*,itemsep=0.2em]
\item The club has stopped serving the Johnsons.
\item If the Johnsons want to, they can serve themselves.
\item The Johnsons think the service is bad.
\item The Johnsons are going to the club.
\end{enumerate}
\par
\vspace{0.25\baselineskip}
\noindent\textbf{18.}
\begin{enumerate}[label=(\Alph*),leftmargin=*,itemsep=0.2em]
\item The new paint makes the house look clean.
\item They worked for three days.
\item They. ran out of paint on the second day.
\item They were in pain for two days.
\end{enumerate}
\par
\vspace{0.25\baselineskip}
\noindent\textbf{19.}
\begin{enumerate}[label=(\Alph*),leftmargin=*,itemsep=0.2em]
\item There are more students in the first class than in the second
\item First class mail delivery is etter than second.
\item First class postage is more expensive than second class.
\item The packaging comes in two prices, first and second.
\end{enumerate}
\par
\vspace{0.25\baselineskip}
\noindent\textbf{20.}
\begin{enumerate}[label=(\Alph*),leftmargin=*,itemsep=0.2em]
\item Mike didn't sleep well.
\item Mike slipped last night.
\item The car is parked for the night.
\item The neighbor has a white dog.
\end{enumerate}
\par
\vspace{0.25\baselineskip}
\noindent\textbf{21.}
\begin{enumerate}[label=(\Alph*),leftmargin=*,itemsep=0.2em]
\item In a taxi
\item In a hotel
\item Downstairs
\item Downtown
\end{enumerate}
\par
\vspace{0.25\baselineskip}
\noindent\textbf{22.}
\begin{enumerate}[label=(\Alph*),leftmargin=*,itemsep=0.2em]
\item The snow will stop soon.
\item The man is wrong about the situation.
\item The man is happy.
\item A call isn't necessary.
\end{enumerate}
\par
\vspace{0.25\baselineskip}
\noindent\textbf{23.}
\begin{enumerate}[label=(\Alph*),leftmargin=*,itemsep=0.2em]
\item Make up her mind
\item Take the man to work
\item Go to school
\item Do the right thing
\end{enumerate}
\par
\vspace{0.25\baselineskip}
\noindent\textbf{24.}
\begin{enumerate}[label=(\Alph*),leftmargin=*,itemsep=0.2em]
\item He doesn't like the company.
\item He has ordered some books.
\item He is a businessman.
\item He is selfish.
\end{enumerate}
\par
\vspace{0.25\baselineskip}
\noindent\textbf{25.}
\begin{enumerate}[label=(\Alph*),leftmargin=*,itemsep=0.2em]
\item A picture
\item A brochure
\item A letter
\item A flier
\end{enumerate}
\par
\vspace{0.25\baselineskip}
\noindent\textbf{26.}
\begin{enumerate}[label=(\Alph*),leftmargin=*,itemsep=0.2em]
\item Joe is not French.
\item Joe is going the wrong way.
\item Joe has passed his exam.
\item Joe will probably fail.
\end{enumerate}
\par
\vspace{0.25\baselineskip}
\noindent\textbf{27.}
\begin{enumerate}[label=(\Alph*),leftmargin=*,itemsep=0.2em]
\item Room 217
\item Room 17
\item The conference room
\item The meeting room
\end{enumerate}
\par
\vspace{0.25\baselineskip}
\noindent\textbf{28.}
\begin{enumerate}[label=(\Alph*),leftmargin=*,itemsep=0.2em]
\item They don't sell milk at the store.
\item They already have enough milk.
\item It's their breakfast.
\item It's too late to go out.
\end{enumerate}
\par
\vspace{0.25\baselineskip}
\noindent\textbf{29.}
\begin{enumerate}[label=(\Alph*),leftmargin=*,itemsep=0.2em]
\item The woman
\item A bakery
\item The woman's husband
\item The woman's mother
\end{enumerate}
\par
\vspace{0.25\baselineskip}
\noindent\textbf{30.}
\begin{enumerate}[label=(\Alph*),leftmargin=*,itemsep=0.2em]
\item The Stevensons have a lot of money.
\item The Stevensons will have to make many arrangements.
\item Will they go by boat or by plane?
\item The Stevensons don't appear to be rich.
\end{enumerate}
\par
\vspace{0.25\baselineskip}
\subsection*{Part B}
DIRECTIONS: In this part of the test, you will hear longer conversations. After each conversation, you will hear several questions. The conversations and questions will not be repeated. After you hear a question, read the four possible answers in. your test book and select. The best answer. Then, on your answer sheet, find the number of the question and fill in the space that corresponds to the letter of the answer you have chosen. Remember, you are not allowed to take notes or write in your test book.\par
\begin{exampleblock}
Listen to the following example:\par
You will hear:\par
You will read:\par
(A) He has changed jobs.\par
(B) He has two children.\par
(C) He has two jobs.\par
(D) He is looking for a job.\par
From the conversation you learn that Tom has taken an additional job. The best answer to the question "Why is Tom tired?" is (C), "He has two jobs." Therefore the correct answer is (C).\par
\end{exampleblock}
\noindent\textbf{31.}
\begin{enumerate}[label=(\Alph*),leftmargin=*,itemsep=0.2em]
\item Leaves
\item Cloth
\item Skins
\item Fibers
\end{enumerate}
\par
\vspace{0.25\baselineskip}
\noindent\textbf{32.}
\begin{enumerate}[label=(\Alph*),leftmargin=*,itemsep=0.2em]
\item Elegant
\item Functional
\item Pretentious
\item Sturdy
\end{enumerate}
\par
\vspace{0.25\baselineskip}
\noindent\textbf{33.}
\begin{enumerate}[label=(\Alph*),leftmargin=*,itemsep=0.2em]
\item 100 years
\item 300 years
\item 600 years
\item 1600 years
\end{enumerate}
\par
\vspace{0.25\baselineskip}
\noindent\textbf{34.}
\begin{enumerate}[label=(\Alph*),leftmargin=*,itemsep=0.2em]
\item Leather has become cheaper.
\item Computers have come down in price.
\item Shoes are made by machines.
\item The cost of labor has declined.
\end{enumerate}
\par
\vspace{0.25\baselineskip}
\noindent\textbf{35.}
\begin{enumerate}[label=(\Alph*),leftmargin=*,itemsep=0.2em]
\item In an airport
\item In a bus depot
\item At a wedding
\item On a school campus
\end{enumerate}
\par
\vspace{0.25\baselineskip}
\noindent\textbf{36.}
\begin{enumerate}[label=(\Alph*),leftmargin=*,itemsep=0.2em]
\item To Baltimore
\item To Philadelphia
\item To Chicago
\item To New York
\end{enumerate}
\par
\vspace{0.25\baselineskip}
\noindent\textbf{37.}
\begin{enumerate}[label=(\Alph*),leftmargin=*,itemsep=0.2em]
\item While visiting their children in college
\item While the man was attending a convention
\item When they went to college together
\item When the woman was in high school
\end{enumerate}
\par
\vspace{0.25\baselineskip}
\noindent\textbf{38.}
\begin{enumerate}[label=(\Alph*),leftmargin=*,itemsep=0.2em]
\item One
\item Two
\item Three
\item Four
\end{enumerate}
\par
\vspace{0.25\baselineskip}
\subsection*{Part C}
DIRECTIONS: In Part C you will hear short lectures and extended conversations. At the end of each, you will be asked several questions. Each lecture or conversation and each question will be spoken only one time. For this reason, you must listen carefully to understand what each speaker says. After you hear a question, read the four choices and select the one that best answers the question the speaker asked. Then, on your answer sheet, find the number of the question and blacken the space that contains the letter for the answer you have chosen. Answer all questions according to what is stated or implied in the lecture or conversation.\par
\begin{exampleblock}
Listen to this sample talk.\par
You will hear:\par
Now listen, to the following example.\par
You will hear:\par
You will read:\par
(A) By cars and carriages\par
(B) By bicycles, trains, and carriages\par
(C) On foot and by boat\par
(D) On board ships and trains\par
The best answer to the question "According to the speaker, how did people travel before the invention of the automobile?" is (B), "By bicycles, trains, and carriages." Therefore, the correct answer is (B). You are not allowed to make notes during the test.\par
\end{exampleblock}
\noindent\textbf{39.}
\begin{enumerate}[label=(\Alph*),leftmargin=*,itemsep=0.2em]
\item How thrilling police work can be
\item Why police work is enlightening
\item That police work has its mundane aspects
\item That serving on the police force is difficult.
\end{enumerate}
\par
\vspace{0.25\baselineskip}
\noindent\textbf{40.}
\begin{enumerate}[label=(\Alph*),leftmargin=*,itemsep=0.2em]
\item As full of risk and adventure
\item As straightforward but demanding
\item As routine but necessary
\item As guaranteed employment
\end{enumerate}
\par
\vspace{0.25\baselineskip}
\noindent\textbf{41.}
\begin{enumerate}[label=(\Alph*),leftmargin=*,itemsep=0.2em]
\item Directing traffic
\item Enforcing speed laws
\item Patrolling streets
\item Resolving disputes
\end{enumerate}
\par
\vspace{0.25\baselineskip}
\noindent\textbf{42.}
\begin{enumerate}[label=(\Alph*),leftmargin=*,itemsep=0.2em]
\item Changing tires for motorists
\item Verifying vehicle licenses
\item Walking the beat on city streets
\item Producing reports and correspondence
\end{enumerate}
\par
\vspace{0.25\baselineskip}
\noindent\textbf{43.}
\begin{enumerate}[label=(\Alph*),leftmargin=*,itemsep=0.2em]
\item It's physically exhausting.
\item It's ineffectual.
\item It can be boring.
\item It may be exhilarating.
\end{enumerate}
\par
\vspace{0.25\baselineskip}
\noindent\textbf{44.}
\begin{enumerate}[label=(\Alph*),leftmargin=*,itemsep=0.2em]
\item The first American government
\item Writing in the early colonial times
\item Population mobility in the colonial times
\item The establishment of first libraries
\end{enumerate}
\par
\vspace{0.25\baselineskip}
\noindent\textbf{45.}
\begin{enumerate}[label=(\Alph*),leftmargin=*,itemsep=0.2em]
\item To describe their new lives
\item To produce literary volumes
\item To conquer wilderness
\item To improve their writing skills
\end{enumerate}
\par
\vspace{0.25\baselineskip}
\noindent\textbf{46.}
\begin{enumerate}[label=(\Alph*),leftmargin=*,itemsep=0.2em]
\item To gain insight and understanding
\item To help shape them into novels
\item To make thema part of legal institutions
\item To develop American national character
\end{enumerate}
\par
\vspace{0.25\baselineskip}
\noindent\textbf{47.}
\begin{enumerate}[label=(\Alph*),leftmargin=*,itemsep=0.2em]
\item Jamestown colony
\item Virginia Company
\item English tailors
\item Religious organizations
\end{enumerate}
\par
\vspace{0.25\baselineskip}
\noindent\textbf{48.}
\begin{enumerate}[label=(\Alph*),leftmargin=*,itemsep=0.2em]
\item In Virginia
\item In Pennsylvania
\item In Boston Bay
\item In Middle Colonies
\end{enumerate}
\par
\vspace{0.25\baselineskip}
\noindent\textbf{49.}
\begin{enumerate}[label=(\Alph*),leftmargin=*,itemsep=0.2em]
\item Philadelphia
\item London
\item Jamestown
\item Salem
\end{enumerate}
\par
\vspace{0.25\baselineskip}
\noindent\textbf{50.}
\begin{enumerate}[label=(\Alph*),leftmargin=*,itemsep=0.2em]
\item Chronicles
\item Observations and essays
\item Government documents
\item Novels and poetry
\end{enumerate}
\par
\vspace{0.25\baselineskip}
\section*{Structure and Written Expression}
\subsection*{Sentence Completion}
This section is designed to test your ability to recognize language structures that are appropriate in standard written English. The questions in this section belong to two types, each of which has special directions.\par
DIRECTIONS: Questions 1-15 are partial sentences. Below each sentence you will see four words or phrases, marked (A), (B), (C), and (D). Select the one word or phrase that best completes the sentence. Then, on your answer sheet, find the number of the question and fill in the space that contains the letter for the answer you have chosen. Fill in the space completely.\par
\begin{exampleblock}
\textbf{EXAMPLE I}\par
\noindent Drying flowers is the best way......them.
\begin{enumerate}[label=(\Alph*),leftmargin=*,itemsep=0.2em]
\item to preserve
\item by preserving
\item preserve
\item preserved
\end{enumerate}
\vspace{0.25\baselineskip}
The sentence should state, "Drying flowers is the best way to preserve them." Therefore, the correct answer is (A).\par
\end{exampleblock}
\begin{exampleblock}
\textbf{EXAMPLE II}\par
\noindent Many American universities......as small, private colleges.
\begin{enumerate}[label=(\Alph*),leftmargin=*,itemsep=0.2em]
\item begun
\item beginning
\item began
\item for the beginning
\end{enumerate}
\vspace{0.25\baselineskip}
The sentence should state, "Many American universities began as small, private colleges." Therefore, the correct answer is (C). After you read the directions, begin work on the questions.\par
\end{exampleblock}
\noindent\textbf{1.} The Boston Public Library,......1854, was the first library to be financed by donations and proceeds from raffles.
\begin{enumerate}[label=(\Alph*),leftmargin=*,itemsep=0.2em]
\item found it in
\item founded in
\item was founded in
\item it was found
\end{enumerate}
\par
\vspace{0.25\baselineskip}
\noindent\textbf{2.} Toolmakers not only...... elaborate tools but also test them for reliability and utility.
\begin{enumerate}[label=(\Alph*),leftmargin=*,itemsep=0.2em]
\item does it help to construct
\item help in the construction
\item help to construct
\item do help to construct
\end{enumerate}
\par
\vspace{0.25\baselineskip}
\noindent\textbf{3.} Before he turned 14, Mozart few lesser pieces for the piano.... a.
\begin{enumerate}[label=(\Alph*),leftmargin=*,itemsep=0.2em]
\item had composed
\item has composed
\item had the composition
\item he had composed
\end{enumerate}
\par
\vspace{0.25\baselineskip}
\noindent\textbf{4.} One of the quickest methods..... personality is the self-report inventory.
\begin{enumerate}[label=(\Alph*),leftmargin=*,itemsep=0.2em]
\item the measuring of
\item to the measures
\item for measurements
\item of measuring
\end{enumerate}
\par
\vspace{0.25\baselineskip}
\noindent\textbf{5.} Ginkgo trees bear seeds.... an unpleasant odor to discourage animals from eating them.
\begin{enumerate}[label=(\Alph*),leftmargin=*,itemsep=0.2em]
\item who have
\item that have
\item which they
\item that are
\end{enumerate}
\par
\vspace{0.25\baselineskip}
\noindent\textbf{6.} Patrick Henry, born in 1736,.... by his father, who had advanced training in mechanics.
\begin{enumerate}[label=(\Alph*),leftmargin=*,itemsep=0.2em]
\item had taught
\item has been taught
\item taught
\item was taught
\end{enumerate}
\par
\vspace{0.25\baselineskip}
\noindent\textbf{7.} In the 1950s, many people believed that the more they produced and consumed,.....
\begin{enumerate}[label=(\Alph*),leftmargin=*,itemsep=0.2em]
\item they were the more affluent
\item the more affluent they were
\item were they affluent
\item they were affluent
\end{enumerate}
\par
\vspace{0.25\baselineskip}
\noindent\textbf{8.} A meteor burns brightly.... through the earth's atmosphere.
\begin{enumerate}[label=(\Alph*),leftmargin=*,itemsep=0.2em]
\item as it descends
\item as descending
\item whether it descends
\item when descends
\end{enumerate}
\par
\vspace{0.25\baselineskip}
\noindent\textbf{9.} Mass transit systems relieve traffic congestion when the service is convenient,....., and affordable.
\begin{enumerate}[label=(\Alph*),leftmargin=*,itemsep=0.2em]
\item comfort
\item comforting
\item comfortable
\item comforted
\end{enumerate}
\par
\vspace{0.25\baselineskip}
\noindent\textbf{10.} Automatic fire alarms,..... smoke detector, are installed in almost all public buildings.
\begin{enumerate}[label=(\Alph*),leftmargin=*,itemsep=0.2em]
\item such as the ubiquitous
\item so ubiquitous
\item such is the ubiquitous
\item so as the ubiquitous
\end{enumerate}
\par
\vspace{0.25\baselineskip}
\noindent\textbf{11.} When a person is in shock, the blood..... enough oxygen to the brain.
\begin{enumerate}[label=(\Alph*),leftmargin=*,itemsep=0.2em]
\item supplies fail to
\item supplied fails to
\item failing to supply
\item fails to supply
\end{enumerate}
\par
\vspace{0.25\baselineskip}
\noindent\textbf{12.} Throughout history, elevated ranges have been viewed as barriers to transportation and.....
\begin{enumerate}[label=(\Alph*),leftmargin=*,itemsep=0.2em]
\item to the communicating
\item to the communications
\item communication
\item communication with
\end{enumerate}
\par
\vspace{0.25\baselineskip}
\noindent\textbf{13.} About 90 percent of fabrics distributed to secondary sewing outlets..... weaving or knitting.
\begin{enumerate}[label=(\Alph*),leftmargin=*,itemsep=0.2em]
\item are manufactured the
\item are manufactured by
\item by manufacturing the
\item manufactured by
\end{enumerate}
\par
\vspace{0.25\baselineskip}
\noindent\textbf{14.} Deer ticks..... vacationers hiking or camping in mixed deciduous forests.
\begin{enumerate}[label=(\Alph*),leftmargin=*,itemsep=0.2em]
\item never trouble
\item never any trouble
\item troubles never
\item trouble never
\end{enumerate}
\par
\vspace{0.25\baselineskip}
\noindent\textbf{15.} Either the goalkeeper or one of the other players.... the ball from the goal.
\begin{enumerate}[label=(\Alph*),leftmargin=*,itemsep=0.2em]
\item retrieving
\item retrieval
\item retrieves
\item retrieve
\end{enumerate}
\par
\vspace{0.25\baselineskip}
\subsection*{Error Identification}
DIRECTIONS: In questions 16-40 every sentence has four words or phrases that are underlined. The four underlined portions of each sentence are marked (A), (B), (C), and (D). Identify the one word or phrase that makes the sentence incorrect. Then, on your answer sheet, find the number of the question and fill in the space that contains the letter for the answer you have chosen. Fill in the space completely.\par
\begin{exampleblock}
\textbf{EXAMPLE I}\par
Christopher Columbus \uline{has sailed}(A) from Europe in 1492 and \uline{discovered a}(B) \uline{new land}(C) he \uline{thought to}(D) be India.\par
The sentence should state, "Christopher Columbus sailed from Europe in 1492 and discovered a new land he thought to be India." Therefore, you should choose answer (A).\par
\end{exampleblock}
\begin{exampleblock}
\textbf{EXAMPLE II}\par
\uline{As the roles}(A) of people in society change, \uline{so does}(B) the \uline{rules of}(C) conduct in certain \uline{situations}(D).\par
The sentence should state, "As the roles of people in society change, so do the rules of\par
conduct in certain situations." Therefore, you should choose answer (B).\par
\end{exampleblock}
\noindent\textbf{16.} The Slater Mill, \uline{built in}(A) 1793, \uline{it was}(B) one of the first \uline{successful}(C) \uline{mills in}(D) the United States.
\par
\vspace{\questionblockgap}
\noindent\textbf{17.} In kindergarten, \uline{children}(A) are generally \uline{unrestricted in}(B) expressing \uline{their ideas}(C) by \uline{talk}(D).
\par
\vspace{\questionblockgap}
\noindent\textbf{18.} \uline{Japanese}(A) initially \uline{used}(B) \uline{jeweled}(C) objects to \uline{decorate}(D) swords and ceremonial items.
\par
\vspace{\questionblockgap}
\noindent\textbf{19.} The legal \uline{age which}(A) a person \uline{is considered}(B) to \uline{be an}(C) adult \uline{is customarily}(D) 18.
\par
\vspace{\questionblockgap}
\noindent\textbf{20.} Australian aborigines \uline{adhere}(A) to their tribal traditions and \uline{few}(B) \uline{marriage}(C) \uline{outside}(D) the tribe.
\par
\vspace{\questionblockgap}
\noindent\textbf{21.} During a radio \uline{broadcast}(A), a microphone \uline{picks up}(B) speech and \uline{another}(C) live \uline{sounds}(D).
\par
\vspace{\questionblockgap}
\noindent\textbf{22.} Although both are the bread and butter of \uline{recreational}(A) vehicles, camping trailers are \uline{smaller}(B) and \uline{compacter}(C) than \uline{travel}(D) trailers.
\par
\vspace{\questionblockgap}
\noindent\textbf{23.} The \uline{leathery fruit}(A) burr of the horse chestnut \uline{splits openly}(B) when \uline{ripe and releases}(C) a \uline{roundish brown}(D) seed.
\par
\vspace{\questionblockgap}
\noindent\textbf{24.} Blacksnakes \uline{ascend}(A) trees to reach \uline{bird's}(B) \uline{nests}(C) and ingest the eggs and young \uline{birds}(D).
\par
\vspace{\questionblockgap}
\noindent\textbf{25.} Although rhubarb is \uline{technically}(A) a vegetable, \uline{it usually}(B) \uline{prepared}(C) as a \uline{dessert}(D).
\par
\vspace{\questionblockgap}
\noindent\textbf{26.} Colonial \uline{craftsmen}(A) \uline{pieced}(B) bed covers together from scraps of linen and wool \uline{because}(C) cloth was \uline{scarcely}(D).
\par
\vspace{\questionblockgap}
\noindent\textbf{27.} In art, \uline{relief}(A) is \uline{sculpture}(B) \uline{in which}(C) the figures or designs \uline{projects}(D) from their background.
\par
\vspace{\questionblockgap}
\noindent\textbf{28.} Edith Roosevelt was a \uline{devoted}(A) mother \uline{of}(B) five children, \uline{as well}(C) a \uline{gracious}(D) hostess.
\par
\vspace{\questionblockgap}
\noindent\textbf{29.} Vocational counseling \uline{guides}(A) students and helps them \uline{to understand}(B) how \uline{occupations differ}(C) and what job opportunities \uline{are exist}(D).
\par
\vspace{\questionblockgap}
\noindent\textbf{30.} Newtonian physics \uline{holds true}(A) if the \uline{velocities}(B) of the objects \uline{being study}(C) are \uline{negligible}(D).
\par
\vspace{\questionblockgap}
\noindent\textbf{31.} Roman doctrine \uline{stipulated}(A) every man \uline{was}(B) born with a \uline{spiritual}(C) who \uline{guarded}(D) him against travail.
\par
\vspace{\questionblockgap}
\noindent\textbf{32.} \uline{When}(A) wine grapes \uline{contain}(B) the proper \uline{amounts}(C) of acid and sugar \uline{required}(D) to produce wine.
\par
\vspace{\questionblockgap}
\noindent\textbf{33.} Beef and dairy cattle \uline{is}(A) major \uline{sources}(B) of income in Louisiana, \uline{which}(C) has a \uline{mild climate}(D).
\par
\vspace{\questionblockgap}
\noindent\textbf{34.} After the new dollar \uline{bills}(A) \uline{are printed}(B) and cut, the inspectors \uline{scrutinize}(C) them for \uline{imperfectives}(D).
\par
\vspace{\questionblockgap}
\noindent\textbf{35.} Psychologists \uline{take it for granted}(A) \uline{that}(B) girls are more \uline{empathetic}(C) than \uline{do}(D) boys.
\par
\vspace{\questionblockgap}
\noindent\textbf{36.} Henry Richardson was the \uline{first}(A) prominent architect to \uline{incorporate}(B) geometric \uline{form}(C) in his \uline{concave designs}(D).
\par
\vspace{\questionblockgap}
\noindent\textbf{37.} With \uline{small numbers}(A), the objects \uline{in a set}(B) can be \uline{visualized}(C) and \uline{quick counted}(D) without mathematical formulas.
\par
\vspace{\questionblockgap}
\noindent\textbf{38.} In \uline{group dancing}(A), couples \uline{step}(B) in tandem, bow, \uline{join hands}(C), and \uline{change partner}(D).
\par
\vspace{\questionblockgap}
\noindent\textbf{39.} Additives are \uline{chemicals}(A) infused into \uline{perishable}(B) foods to prevent \uline{it}(C) from \uline{spoiling}(D).
\par
\vspace{\questionblockgap}
\noindent\textbf{40.} Football is a \uline{fast-moving}(A) \uline{team sport}(B) \uline{playing}(C) \uline{mainly}(D) in the United States and Canada.
\par
\vspace{\questionblockgap}
\section*{Reading Comprehension}
\begin{center}
\textbf{Time 55 minutes}
\end{center}
DIRECTIONS: In this section you will read several passages. Each passage is followed by a series of questions. For questions 1-50, you need to select the best answer, (A), (B), (C), or (D), to each question. Then, on your answer sheet, find the number of the question and fill in the space that contains the letter of the answer you have selected. Fill in the space completely. Answer all questions following a passage on the basis of what is stated or implied in the passage.\par
\begin{exampleblock}
\small
\sloppy
Read the following passage:\par
A tomahawk is a small ax used as a tool and a weapon by the North American Indian\par
tribes. An average tomahawk was not very long and did not weigh a great deal. Originally,\par
the head of the tomahawk was made of a shaped stone or an animal bone and was mounted on\par
Line a wooden handle. After the arrival of the European settlers, the Indians began to use toma-\par
hawks with iron heads. Indian males and females of all ages used tomahawks to chop and cut wood, pound stakes into the ground to put up wigwams, and do many other chores. Indian warriors relied on tomahawks as weapons and even threw them at their enemies. Some types of tomahawks were used in religious ceremonies. Contemporary American idioms reflect\par
this aspect of American heritage.\par
\end{exampleblock}
\begin{exampleblock}
\textbf{EXAMPLE I}\par
\noindent Early tomahawk heads were made of
\begin{enumerate}[label=(\Alph*),leftmargin=*,itemsep=0.2em]
\item stone or bone
\item wood or sticks
\item European iron
\item religious weapons
\end{enumerate}
\vspace{0.25\baselineskip}
According to the passage, early tomahawk heads were made of stone or bone. Therefore, the correct answer is (A).\par
\end{exampleblock}
\begin{exampleblock}
\textbf{EXAMPLE II}\par
\noindent How has the Indian use of tomahawks affected American daily life today?
\begin{enumerate}[label=(\Alph*),leftmargin=*,itemsep=0.2em]
\item Tomahawks are still used as weapons.
\item Tomahawks are used as tools for.certain jobs.
\item Contemporary language refers to tomahawks.
\item Indian tribes cherish tomahawks as heirlooms.
\end{enumerate}
\vspace{0.25\baselineskip}
The passage states, "Contemporary American idioms reflect this aspect of American heritage." The correct answer is (C). After you read the directions, begin work on the questions.\par
\end{exampleblock}
\subsection*{Questions 1-11}
\begingroup
\small
\sloppy
\setlength{\parskip}{0pt}
\setlength{\parindent}{0pt}
The Mayo Clinic in Rochester, Minnesota, where staff physicians practice a special\par
integrated approach to patient care, is one of the largest medical facilities in the world. The\par
clinic staff consists of a 12-member, committee-based board of governors and 900 physicians and medical personnel whose records are updated by approximately 200 auxiliary per-\par
sonnel. About 800 resident doctors assist the full-time physicians as a phase of their training in medicine and surgery while they acquire their specializations. The Mayo approach to\par
treatment has been hailed for its almost miraculous patient recovery rate.\par
William Worral Mayo was born in Manchester, England, immigrated to the United\par
States in 1845, and immediately began his medical training. In 1860, he took an active part\par
in organizing the Minnesota Territory and accepted the position of an Army surgeon during\par
a Sioux Indian outbreak. This appointment became a stepping stone for his advancement\par
to the post of provost surgeon for the southern portion of the state in 1863. His personal\par
dedication and courage became legendary when a cyclone struck Rochester, and he was\par
placed in charge of an emergency hospital.\par
William Worral Mayo provided crucial assistance to his sons in launching their team\par
practice in 1889, while they were holding positions at St. Mary's Hospital. William James\par
became recognized for his surgical skill in gallstone, cancer, and abdominal operations. He\par
and his brother, Charles Horace, founded the Mayo Graduate School of Medicine and donated \$1.5 million to establish the foundation for contributions. Charles William, the son(20) of Charles Horace Mayo, became a member of the board of governors at the Mayo Graduate\par
School and an alternate delegate to the United Nations General Assembly before retiring\par
from the clinic in 1963.\par
William James Mayo presided in the American Medical Association and served in the\par
Army military corps as a brigadier general in the medical reserve. Charles Horace was a pro-\par
fessor of surgery and a health officer of Rochester subsequent to serving in the armed forces\par
between 1914 and 1918. The Mayo practice became known far and wide for its success in\par
surgical procedures. In 1914, the practice moved into its.own medical center, and today the\par
number of patients equals approximately 280,000 per annum. Since the clinic opened in\par
1907, 4.5 million patients have been treated there..\par
\endgroup
\medskip
\noindent\textbf{1.} What is the best title for the passage?
\begin{enumerate}[label=(\Alph*),leftmargin=*,itemsep=0.2em]
\item A Welcome to the Mayo Clinic
\item Brilliant American Surgeons
\item The Contributions of the Mayo Family
\item The Start of a Successful Practice
\end{enumerate}
\par
\vspace{0.25\baselineskip}
\noindent\textbf{2.} In line 2, the word "integrated" is closest in meaning to
\begin{enumerate}[label=(\Alph*),leftmargin=*,itemsep=0.2em]
\item unified
\item goal-oriented
\item ready-made
\item unique
\end{enumerate}
\par
\vspace{0.25\baselineskip}
\noindent\textbf{3.} What can be inferred from the first paragraph?
\begin{enumerate}[label=(\Alph*),leftmargin=*,itemsep=0.2em]
\item Rochester, Minnesota is a city with a large population.
\item The Mayo Clinic employees have set many medical records.
\item The Mayo Clinic is a large specialized teaching hospital.
\item The clinic's physicians represent many medical specializations..
\end{enumerate}
\par
\vspace{0.25\baselineskip}
\noindent\textbf{4.} According to the passage, William Worral Mayo was involved in caring for the patients affected by
\begin{enumerate}[label=(\Alph*),leftmargin=*,itemsep=0.2em]
\item an outbreak of an epidemic
\item a spread of disease among Indians
\item a devastating natural disaster
\item a brief military confrontation
\end{enumerate}
\par
\vspace{0.25\baselineskip}
\noindent\textbf{5.} In line 11, the word "dedication" is closest in meaning to
\begin{enumerate}[label=(\Alph*),leftmargin=*,itemsep=0.2em]
\item dejection
\item devotion
\item deliberation
\item delectation
\end{enumerate}
\par
\vspace{0.25\baselineskip}
\noindent\textbf{6.} According to the passage, who were the first physicians in the clinic?
\begin{enumerate}[label=(\Alph*),leftmargin=*,itemsep=0.2em]
\item William Worral Mayo and Charles Horace Mayo
\item William James Mayo and Charles Horace Mayo
\item William Worral Mayo and William James Mayo
\item William James Mayo and Charles William Mayo
\end{enumerate}
\par
\vspace{0.25\baselineskip}
\noindent\textbf{7.} In line 16, the word "founded" is closest in meaning to
\begin{enumerate}[label=(\Alph*),leftmargin=*,itemsep=0.2em]
\item found
\item established
\item fortified
\item articulated
\end{enumerate}
\par
\vspace{0.25\baselineskip}
\noindent\textbf{8.} In line 16, the word "contributions" is closest in meaning to
\begin{enumerate}[label=(\Alph*),leftmargin=*,itemsep=0.2em]
\item contractions
\item conventions
\item gifts
\item prizes
\end{enumerate}
\par
\vspace{0.25\baselineskip}
\noindent\textbf{9.} In addition to their medical expertise, what common characteristics distinguished the careers of the Mayo brothers?
\begin{enumerate}[label=(\Alph*),leftmargin=*,itemsep=0.2em]
\item Funding of schools in Minnesota
\item Positions on the board of directors
\item Military and political service
\item Donations for poor patients
\end{enumerate}
\par
\vspace{0.25\baselineskip}
\noindent\textbf{10.} In line 21, the phrase "subsequent to" is closest in meaning to
\begin{enumerate}[label=(\Alph*),leftmargin=*,itemsep=0.2em]
\item because
\item regardless of
\item after
\item contrary to
\end{enumerate}
\par
\vspace{0.25\baselineskip}
\noindent\textbf{11.} Where in the passage does the author state the principal reason for the expansion of the practice?
\begin{enumerate}[label=(\Alph*),leftmargin=*,itemsep=0.2em]
\item Lines 5-6
\item Lines 11-14
\item Lines 17-22
\item Lines 22-25
\end{enumerate}
\par
\vspace{0.25\baselineskip}
\subsection*{Questions 12-23}
\begingroup
\small
\sloppy
\setlength{\parskip}{0pt}
\setlength{\parindent}{0pt}
Consumers are frequently unaware that about 30 percent of nationwide department\par
stores are franchises with numerous outlets. Chain stores are a group of retail stores that are\par
supervised or coordinated by centralized management. From a business perspective, chain\par
stores have numerous advantages over independent stores, one of which is that the parent\par
company almost always has the credit to purchase large quantities of goods to supply to its\par
outlets and to receive a discount for placing such an order. Through the centralized system\par
of distribution, chain stores can absorb the cost and price differential and attract consumers\par
with various physical and psychological needs. They can also distribute their operating\par
costs for accounting, advertising, marketing, merchandising, and transportation.\par
In general, approximately 50 percent of gross product cost results from the associated\par
marketing research and distribution. While research focuses on the probable market segments, it strongly considers consumer behavior and cognitive motives rather than the actual prices of goods. Similarly, the cost increase in the multiple channels of distribution\par
accounts for about 23 percent of the unit price. By combining their marketing resources and\par
distribution networks, franchise outlets can avoid performing whole stages of marketing\par
studies and layers of distribution networks to reduce unit prices. It is the central company\par
that conducts marketing and communicates with manufacturers, thus controlling production decisions and the pricing policy. Franchises operate according to their contracts with\par
the parent company and pay it a fraction of their net gains. They symbolize a brand name\par
and identify their goods with a particular range of quality that sets it apart from other, similar products. Essentially, chain stores convert consumer brand name loyalty into profit; this\par
factor determines franchise proliferation and results in a relatively low degree of failure.\par
\endgroup
\medskip
\noindent\textbf{12.} With what topic is the passage mainly concerned?
\begin{enumerate}[label=(\Alph*),leftmargin=*,itemsep=0.2em]
\item The marketing of chain store products
\item The business rationale for chain stores
\item Pricing and distribution in franchises
\item Brand name'imaging of retail outlets
\end{enumerate}
\par
\vspace{0.25\baselineskip}
\noindent\textbf{13.} In line 2, the word "outlets" is closest in meaning to
\begin{enumerate}[label=(\Alph*),leftmargin=*,itemsep=0.2em]
\item stock markets
\item store rooms
\item retailers
\item outfits
\end{enumerate}
\par
\vspace{0.25\baselineskip}
\noindent\textbf{14.} In line 7, the word "absorb" is closest in meaning to
\begin{enumerate}[label=(\Alph*),leftmargin=*,itemsep=0.2em]
\item absolve
\item cushion
\item hide
\item advertis
\end{enumerate}
\par
\vspace{0.25\baselineskip}
\noindent\textbf{15.} It can be inferred from the passage that chain stores
\begin{enumerate}[label=(\Alph*),leftmargin=*,itemsep=0.2em]
\item are more expensive than department stores
\item can economize by controlling their operating costs
\item have a greater consumer appeal than the parent company
\item have power in identifying their market segments
\end{enumerate}
\par
\vspace{0.25\baselineskip}
\noindent\textbf{16.} According to the passage, what does marketing research include?
\begin{enumerate}[label=(\Alph*),leftmargin=*,itemsep=0.2em]
\item The actual prices of goods
\item Consumer segments and behavior
\item Multiple channels of distribution
\item The percentage of gross product cost
\end{enumerate}
\par
\vspace{0.25\baselineskip}
\noindent\textbf{17.} In line 12, the phrase "accounts for" is closest in meaning to
\begin{enumerate}[label=(\Alph*),leftmargin=*,itemsep=0.2em]
\item sees as
\item counts on
\item . adjusts
\item represents
\end{enumerate}
\par
\vspace{0.25\baselineskip}
\noindent\textbf{18.} Which of the following is NOT mentioned as a means through which chain stores control their prices?
\begin{enumerate}[label=(\Alph*),leftmargin=*,itemsep=0.2em]
\item Reducing distribution costs
\item Consolidating their finances
\item Dividing their purchase orders
\item Marketing a company brand name
\end{enumerate}
\par
\vspace{0.25\baselineskip}
\noindent\textbf{19.} Why does the author mention legally binding agreements between businesses?
\begin{enumerate}[label=(\Alph*),leftmargin=*,itemsep=0.2em]
\item To show feasible profitability of merchandising
\item To measure the value of consumer dependence on a product
\item To point out the means of parent company control
\item To exemplify the system of franchise Operations
\end{enumerate}
\par
\vspace{0.25\baselineskip}
\noindent\textbf{20.} In line 16, the word "they" refers to
\begin{enumerate}[label=(\Alph*),leftmargin=*,itemsep=0.2em]
\item production decisions and the pricing
\item franchises
\item contracts
\item their net gains
\end{enumerate}
\par
\vspace{0.25\baselineskip}
\noindent\textbf{21.} It can be inferred from the passage that the parent company probably dictates
\begin{enumerate}[label=(\Alph*),leftmargin=*,itemsep=0.2em]
\item what sales personnel are employed
\item what profit an outlet makes
\item how goods are advertised
\item how products are packed
\end{enumerate}
\par
\vspace{0.25\baselineskip}
\noindent\textbf{22.} According to the passage, how do chain stores profit by being associated with the parent company?
\begin{enumerate}[label=(\Alph*),leftmargin=*,itemsep=0.2em]
\item They are not concerned about market instability,
\item They are not held accountable for a change in profits.
\item They put their needs ahead of those of the parent company.
\item They market consumer brand name recognition.
\end{enumerate}
\par
\vspace{0.25\baselineskip}
\noindent\textbf{23.} In line 17, the word "convert" is closest in meaning to
\begin{enumerate}[label=(\Alph*),leftmargin=*,itemsep=0.2em]
\item turn
\item nestle
\item antagonize
\item fit
\end{enumerate}
\par
\vspace{0.25\baselineskip}
\subsection*{Questions 24-35}
\begingroup
\small
\sloppy
\setlength{\parskip}{0pt}
\setlength{\parindent}{0pt}
Because conducting censuses requires detailed planning, the organization conducting a\par
census decides on and narrows the topics to be addressed and, more specifically, determines\par
how to word the questions, tabulate the responses, and report the findings. Assimilating,\par
compiling, and statistically analyzing the information is a work-intensive process that may\par
sometimes take up to a year. Censuses examine such issues as population size and density,\par
employment and industrial affiliation, migration, formal education, income received, housing, marital status, relationship of each individual to the head of the family, and age. A detailed series of queries sample the data associated with the quality of housing, transportation, the level of industrial production, water and electricity consumption, or the\par
functioning of the local government. Major censuses taken by the federal government are\par
conducted every 10 years, in the years that end in zero. Surveys of agriculture take place\par
every 5 years and cover the years ending in 4 and 9, manufacturing censuses in the years that\par
end in 3 and 8, governmental units in the years ending in 2 and 7, and drainage and irrigation systems in the years ending in 9.\par
To ensure that the census information is complete, the organization conducting the census attempts to contact every individual residing or employed in a particular geographic\par
area. To be consistent, information is gathered at approximately the same time. Because it is\par
practically impossible to reach every person on the same day, censuses question the individual about conditions as they were on a certain date. Following the data gathering, the infor-\par
mation is analyzed to determine the extent of social and economic change and problems, as\par
well as the resources available to deal with them. During the years between censuses, the\par
Census Bureau engages in monthly interviews and queries from a sample population to update its statistics. The issues of Statistical Abstract of the United States summarize all the information that is collected by 50 federal agencies and by private and public agencies.\par
\endgroup
\medskip
\noindent\textbf{24.} With what aspect of conducting censuses is the passage mainly concerned?
\begin{enumerate}[label=(\Alph*),leftmargin=*,itemsep=0.2em]
\item Content and timing
\item Inquiring and analyzing
\item Tallying and updating
\item Reporting and publicizing
\end{enumerate}
\par
\vspace{0.25\baselineskip}
\noindent\textbf{25.} In line 2, the word "narrows" is closest in meaning to
\begin{enumerate}[label=(\Alph*),leftmargin=*,itemsep=0.2em]
\item delimits
\item declines
\item defrays
\item deflects
\end{enumerate}
\par
\vspace{0.25\baselineskip}
\noindent\textbf{26.} According to the passage, censuses take a great deal of
\begin{enumerate}[label=(\Alph*),leftmargin=*,itemsep=0.2em]
\item funding
\item publicity
\item insight
\item time
\end{enumerate}
\par
\vspace{0.25\baselineskip}
\noindent\textbf{27.} The author mentions that censuses gather data about all of the following EXCEPT
\begin{enumerate}[label=(\Alph*),leftmargin=*,itemsep=0.2em]
\item household size
\item agricultural production
\item industrial output
\item social networks
\end{enumerate}
\par
\vspace{0.25\baselineskip}
\noindent\textbf{28.} In line 6, the word "queries" is closest in meaning to
\begin{enumerate}[label=(\Alph*),leftmargin=*,itemsep=0.2em]
\item entities
\item terms
\item quandaries
\item questions
\end{enumerate}
\par
\vspace{0.25\baselineskip}
\noindent\textbf{29.} It can be inferred from the passage that a census of industries took place in
\begin{enumerate}[label=(\Alph*),leftmargin=*,itemsep=0.2em]
\item 1990
\item 1994
\item 1998
\item 1997
\end{enumerate}
\par
\vspace{0.25\baselineskip}
\noindent\textbf{30.} In line 12, the word "consistent" is closest in meaning to
\begin{enumerate}[label=(\Alph*),leftmargin=*,itemsep=0.2em]
\item constant
\item uniform
\item assorted
\item conjectural
\end{enumerate}
\par
\vspace{0.25\baselineskip}
\noindent\textbf{31.} To gather complete information, the organization conducting the census contacts all people who
\begin{enumerate}[label=(\Alph*),leftmargin=*,itemsep=0.2em]
\item speak at the same time
\item have a similar economic status
\item are involved in similar activities
\item are concerned about a specific problem
\end{enumerate}
\par
\vspace{0.25\baselineskip}
\noindent\textbf{32.} Where in the passage does the author describe how census information is collected?
\begin{enumerate}[label=(\Alph*),leftmargin=*,itemsep=0.2em]
\item Lines 1-3
\item Lines 11-13
\item Lines 14-16
\item Lines 7-10
\end{enumerate}
\par
\vspace{0.25\baselineskip}
\noindent\textbf{33.} It can be inferred from the passage that census data are necessary to
\begin{enumerate}[label=(\Alph*),leftmargin=*,itemsep=0.2em]
\item rebuild urban infrastructure
\item determine arising needs
\item analyze the electoral outlook
\item identify sources of criminal activity
\end{enumerate}
\par
\vspace{0.25\baselineskip}
\noindent\textbf{34.} The author of the passage implies that most censuses are conducted by
\begin{enumerate}[label=(\Alph*),leftmargin=*,itemsep=0.2em]
\item private agencies
\item public organizations
\item the central government
\item the statistical bureau
\end{enumerate}
\par
\vspace{0.25\baselineskip}
\noindent\textbf{35.} In line 16, the word "statistics" is closest in meaning to
\begin{enumerate}[label=(\Alph*),leftmargin=*,itemsep=0.2em]
\item equations
\item solutions
\item questionnaires
\item data
\end{enumerate}
\par
\vspace{0.25\baselineskip}
\subsection*{Questions 36-42}
\begingroup
\small
\sloppy
\setlength{\parskip}{0pt}
\setlength{\parindent}{0pt}
The body of the honey bee, like the bodies of all insects, is divided into three sections: the\par
head, the thorax, and the abdomen. The bee's entire body is covered with fine hairs to which\par
grains of pollen adhere as the bee moves from flower to flower, harvesting nectar and pollinating plants. The hairs on the antennae provide a means for tactile sensing without the\par
ability. to grasp extending objects. In pigmentation, bees range from black to shades of very\par
pale-brown, thus reflecting their wasplike ancestry. The queens are larger by far than both\par
workers and drones, with the drones being bigger than workers.\par
A honey bee has five eyes-three secondary ones that form a triangle on top of its head\par
and a large compound eye on either side of its head. The compound eyes center around thou-\par
sands of lenses clustered closely to one another. Bees cannot focus their eyes, as many mammals do, because their eyes have no pupils. Bees were the first insects known to distinguish\par
color, an ability due to the color sensitivity of their optic nerve particles. Their vision is especially receptive to hues of blue and yellow and to ultraviolet rays, unseen by humans. However, bees see red in the same way they see green but can distinguish geometrical patterns in\par
the shapes of foliage and blossoms.\par
\endgroup
\medskip
\noindent\textbf{36.} It can be inferred from the passage that the body of the bee
\begin{enumerate}[label=(\Alph*),leftmargin=*,itemsep=0.2em]
\item has a pronounced abdomen
\item has fine layers
\item is colored for protection
\item is typically compound
\end{enumerate}
\par
\vspace{0.25\baselineskip}
\noindent\textbf{37.} According to the passage, what do bees collect?
\begin{enumerate}[label=(\Alph*),leftmargin=*,itemsep=0.2em]
\item flowers
\item grains
\item nectar and pollen
\item extending objects
\end{enumerate}
\par
\vspace{0.25\baselineskip}
\noindent\textbf{38.} According to the passage, what purpose do the fine hairs on the body of the bee serve?
\begin{enumerate}[label=(\Alph*),leftmargin=*,itemsep=0.2em]
\item They identify the bees' ancestry.
\item They furnish a sense of touch.
\item They carry pollen to female
\item They camouflage the insect.
\end{enumerate}
\par
\vspace{0.25\baselineskip}
\noindent\textbf{39.} How many classes of bees are mentioned in the passage?
\begin{enumerate}[label=(\Alph*),leftmargin=*,itemsep=0.2em]
\item Two
\item Three
\item Four
\item Five
\end{enumerate}
\par
\vspace{0.25\baselineskip}
\noindent\textbf{40.} In line 9, the word "clustered" is closest in meaning to
\begin{enumerate}[label=(\Alph*),leftmargin=*,itemsep=0.2em]
\item lined up
\item turned up
\item bunched
\item bundled
\end{enumerate}
\par
\vspace{0.25\baselineskip}
\noindent\textbf{41.} In line 10, the word "hues" is closest in meaning to
\begin{enumerate}[label=(\Alph*),leftmargin=*,itemsep=0.2em]
\item shades
\item layers
\item specks
\item circles
\end{enumerate}
\par
\vspace{0.25\baselineskip}
\noindent\textbf{42.} It can be inferred from the passage that bees are LEAST likely to distinguish
\begin{enumerate}[label=(\Alph*),leftmargin=*,itemsep=0.2em]
\item ultraviolet light
\item red flowers from foliage
\item shapes in their proximity
\item patterns of leaf veins
\end{enumerate}
\par
\vspace{0.25\baselineskip}
\subsection*{Questions 43-50}
\begingroup
\small
\sloppy
\setlength{\parskip}{0pt}
\setlength{\parindent}{0pt}
The term "fixed star" refers to stars whose positions do not seem to change relative to\par
other stars. As stars are constantly moving, their patterns alter gradually, but the difference\par
may not appear significant over a period of 50 or 70 years, because, unlike planets, stars are\par
located at a great distance from the earth. Barnard's star, which is considered to move the\par
fastest among the fixed stars, requires 200 years to noticeably alter its position by a distance\par
equal to the moon's diameter. Compared to planets, which seem to be continuously shifting, stars appear stationary in their positions and relative distances from other bodies in a\par
constellation. Astronomers employ spectrographs and photometers to keep track of perpetually shifting constellations. By using the spectrographs and photographs obtained over\par
decades, scientists are able to detect changes in the proper motions of stars. By computing\par
the differences in the spectra length, color, and shade of color, they can predict the direction\par
of a star's movement in the near and distant future. A classical example is the cup of the Big\par
Dipper, which is open quite a bit more widely now than it was 40,000 years ago. According\par
to the spectrograph-based projections, 50,000 years from now the opening will be so expan-\par
sive that the constellation will no longer resemble a dipper.\par
\endgroup
\medskip
\noindent\textbf{43.} This passage is probably taken from an article discussing
\begin{enumerate}[label=(\Alph*),leftmargin=*,itemsep=0.2em]
\item the uses of spectra for projections
\item various fixed constellations and stars
\item the classification of planets and stars
\item the emergence and motions of stars
\end{enumerate}
\par
\vspace{0.25\baselineskip}
\noindent\textbf{44.} The author of the passage implies that stars
\begin{enumerate}[label=(\Alph*),leftmargin=*,itemsep=0.2em]
\item seem to rotate differently than planets
\item appear to retain their distances from other stars
\item are presumed to be perpetually fixed and immobile
\item alter their ray emissions gradually
\end{enumerate}
\par
\vspace{0.25\baselineskip}
\noindent\textbf{45.} The author mentions Barnard's star as
\begin{enumerate}[label=(\Alph*),leftmargin=*,itemsep=0.2em]
\item an example of a very brilliant star
\item a description of a constellation divide
\item a point of reference for star mobility
\item proof of star diameter and proximity
\end{enumerate}
\par
\vspace{0.25\baselineskip}
\noindent\textbf{46.} In line 6, the word "stationary" is closest in meaning to
\begin{enumerate}[label=(\Alph*),leftmargin=*,itemsep=0.2em]
\item shining
\item glimmering
\item immobile
\item immense
\end{enumerate}
\par
\vspace{0.25\baselineskip}
\noindent\textbf{47.} It can be inferred from the passage that
\begin{enumerate}[label=(\Alph*),leftmargin=*,itemsep=0.2em]
\item stars revolve every 40,000 to 50,000 years
\item stars alter their trajectories every 200 years
\item planets are closer to the earth than stars
\item planets do not appear to be as close as stars
\end{enumerate}
\par
\vspace{0.25\baselineskip}
\noindent\textbf{48.} In line 8, the word "detect" is closest in meaning to
\begin{enumerate}[label=(\Alph*),leftmargin=*,itemsep=0.2em]
\item spy
\item project
\item trace
\item compute
\end{enumerate}
\par
\vspace{0.25\baselineskip}
\noindent\textbf{49.} The author of the passage implies that astronomers detect the movement of stars by
\begin{enumerate}[label=(\Alph*),leftmargin=*,itemsep=0.2em]
\item using photographs and mathematical reasoning
\item comparing the spectrum of each categorize. constellation
\item overlaying the spectrum images in succession
\item forecasting changes in the position of stars
\end{enumerate}
\par
\vspace{0.25\baselineskip}
\noindent\textbf{50.} The author of the passage conveys which of the following about fixed stars?
\begin{enumerate}[label=(\Alph*),leftmargin=*,itemsep=0.2em]
\item They are inspected by means of photometers.
\item They are comparatively simple to categorize.
\item Their movements are imperceptible.
\item Their spectra vary in intensity.
\end{enumerate}
\par
\vspace{0.25\baselineskip}
\end{document}
